% Native hgraph authoring. Original statement blocks are preserved verbatim.
\section{Framed manifolds of Kervaire invariant one}
% AimPaper: 1.1 theorem 102 104
\begin{theorem}[A framed manifold of Kervaire invariant one in dimension $126$]
\label{paper:1.1}
\source{AimPaper Theorem 1.1}
\level{coarse}
From \cite{LWX126}, Theorem 1.1.\par
% BEGIN AIMPAPER STATEMENT
\label{thm:main}
    There exist framed manifolds of Kervaire invariant one in dimension 126.
% END AIMPAPER STATEMENT
\uses{paper:1.3,paper:1.4}
% Blueprint note: Mathematical dependency review: [AimPaper source, lines 113-123](paper/main.tex.html#L113). Browder's geometric criterion applied to the permanence theorem gives the framed 126-dimensional manifold.

% Correspondence: public correspondence; local deduction proved, geometric upstream obligations pending.
% Blueprint note: The geometric existence predicate retains an actual closed smooth stably framed manifold and its invariant. The local deduction is complete; the invariant construction and Browder comparison remain separately registered Literature obligations.
\lean{KIP126.Geometry.kervaire_invariant_one_126}

\end{theorem}
\begin{proof}
\textit{(See \cite{LWX126}, Theorem 1.1.)} \href{paper/main.tex.html#L102}{Original statement and context}.

The permanence theorem gives nonzero survival of the specified standard square $h_6^2$ in the fixed sphere Adams spectral sequence. Apply Browder's criterion, Theorem~\ref{paper:1.3}, with $j=6$. Its dimension is $2^{6+1}-2=126$, so it yields a closed smooth stably framed manifold in that dimension with Kervaire invariant one. The local deduction uses the geometric invariant and comparison supplied by Literature; those upstream constructions retain their registered pending status.
\end{proof}


% AimPaper: 1.2 corollary 108 110
\begin{corollary}[Dimensions admitting Kervaire invariant one]
\label{paper:1.2}
\source{AimPaper Corollary 1.2}
\level{coarse}
From \cite{LWX126}, Corollary 1.2.\par
% BEGIN AIMPAPER STATEMENT
The dimensions that there exist framed manifolds of Kervaire invariant one are 2, 6, 14, 30, 62, and 126.
% END AIMPAPER STATEMENT
\uses{paper:1.1,paper:1.3,sem:Main:two-line-known-squares,sem:Literature:two-line-classical-inputs}
% Blueprint note: Mathematical dependency review: [AimPaper source, lines 106-114](paper/main.tex.html#L106). Combine the new dimension-126 case with Browder, the known lower Kervaire squares and HHR's exclusion; historical existence results remain bibliography inputs.

% Correspondence: public correspondence; local deduction proved, geometric upstream obligations pending.
% Blueprint note: The theorem proves both directions of the exact dimension list from Browder, the known low squares, the new h6-square theorem and HHR. Its local proof has no placeholders; all transitive upstream responsibilities are retained.
\lean{KIP126.Geometry.kervaire_invariant_one_iff}

\end{corollary}
\begin{proof}
\textit{(See \cite{LWX126}, Corollary 1.2.)} \href{paper/main.tex.html#L108}{Original statement and context}.

By Browder's criterion, Theorem~\ref{paper:1.3}, existence is equivalent to $n=2^{j+1}-2$ and nonzero survival of $h_j^2$, with $j\ge1$. The low squares for $1\le j\le3$ survive by the classical inputs in \cite{MayThesis1964}; $j=4,5$ are supplied by \cite{MahowaldTangora,BJMtheta5}. Theorem~\ref{paper:1.4} supplies $j=6$, while \cite{HHR} excludes $j\ge7$.

Thus precisely $1\le j\le6$ occur. Evaluating $2^{j+1}-2$ at these six indices gives $2,6,14,30,62,126$. Conversely every listed dimension is supplied by its surviving square and Browder's criterion. This requires neither an order-two $\theta_6$ nor a vanishing assertion for $\theta_5^2$; it inherits the same geometric and literature comparison inputs as Theorem~\ref{paper:1.1}.
\end{proof}
